% Standalone problem document, generated from the adjacent README.md.
% Compile with XeLaTeX. Mathematical content is maintained in Markdown.
\documentclass[11pt,a4paper]{article}
\usepackage[fontsize=10.5pt]{scrextend}
\usepackage[margin=22mm,headheight=15pt,headsep=9mm,footskip=11mm]{geometry}
\usepackage{fontspec}
\setmainfont{texgyrepagella}[Extension=.otf,UprightFont=*-regular,BoldFont=*-bold,ItalicFont=*-italic,BoldItalicFont=*-bolditalic]
\setsansfont{texgyreheros}[Extension=.otf,UprightFont=*-regular,BoldFont=*-bold,ItalicFont=*-italic,BoldItalicFont=*-bolditalic]
\setmonofont{texgyrecursor}[Extension=.otf,UprightFont=*-regular,BoldFont=*-bold,ItalicFont=*-italic,BoldItalicFont=*-bolditalic,Scale=MatchLowercase]
\usepackage{amsmath,amssymb}
\usepackage{unicode-math}
\setmathfont{texgyrepagella-math.otf}
\usepackage{xcolor,microtype,enumitem,needspace,fancyhdr}
\usepackage{bookmark}
\definecolor{ink}{HTML}{17344B}
\definecolor{accent}{HTML}{197D80}
\definecolor{muted}{HTML}{596773}
\hypersetup{colorlinks=true,linkcolor=accent,urlcolor=accent,pdftitle={AC-13 - Deterministic
near-quadratic matrix multiplication
verification},pdfauthor={Open Problems in Numerical Linear Algebra}}
\urlstyle{same}
\setlength{\parindent}{0pt}
\setlength{\parskip}{5pt plus 1pt minus 1pt}
\linespread{1.04}
\setlength{\emergencystretch}{3em}
\widowpenalty=10000
\clubpenalty=10000
\displaywidowpenalty=10000
\allowdisplaybreaks[1]
\setlist{nosep,leftmargin=1.5em}
\providecommand{\tightlist}{\setlength{\itemsep}{0pt}\setlength{\parskip}{0pt}}
\providecommand{\pandocbounded}[1]{#1}
\pagestyle{fancy}
\fancyhf{}
\fancyhead[L]{\sffamily\footnotesize\color{muted}OPEN PROBLEMS IN NUMERICAL LINEAR ALGEBRA}
\fancyhead[R]{\sffamily\bfseries\footnotesize\color{accent}AC-13}
\fancyfoot[L]{\parbox[b]{0.9\textwidth}{\sffamily\fontsize{8}{10}\selectfont\color{muted}Editorial ratings; a bounded search cannot certify that a problem remains unsolved.\\Literature check: 2026-09-10}}
\fancyfoot[R]{\sffamily\footnotesize\color{muted}\thepage}
\renewcommand{\headrulewidth}{0.3pt}
\renewcommand{\headrule}{\hbox to\headwidth{\color{accent}\leaders\hrule height \headrulewidth\hfill}}
\makeatletter
\renewcommand{\section}{\Needspace{5\baselineskip}\@startsection{section}{1}{0pt}{12pt plus 2pt minus 2pt}{4pt}{\sffamily\bfseries\large\color{ink}}}
\renewcommand{\subsection}{\Needspace{4\baselineskip}\@startsection{subsection}{2}{0pt}{9pt plus 1pt minus 1pt}{3pt}{\sffamily\bfseries\normalsize\color{ink}}}
\makeatother
\setcounter{secnumdepth}{0}
\begin{document}
{\sffamily\small\color{muted}Arithmetic and complexity\par}
\vspace{3pt}
{\raggedright\hyphenpenalty=10000\exhyphenpenalty=10000\sffamily\bfseries\fontsize{21}{24}\selectfont\color{ink}Deterministic
near-quadratic matrix multiplication verification\par}
\vspace{7pt}
{\sffamily\small\textbf{Difficulty:} extreme\quad\textbf{Importance:} broadly
interesting\par}
{\sffamily\small\bfseries\color{accent}Status: Partially resolved\par}
\vspace{4pt}
{\color{accent}\hrule height 0.8pt}
\vspace{6pt}
\textbf{Rating rationale:} Matching randomized verification in
deterministic near-quadratic bit time is a longstanding derandomization
barrier, with obstructions known for restricted algorithm classes. Exact
matrix verification has broad importance across linear algebra and
computational complexity.

\textbf{Area:} computational complexity; exact linear algebra

\subsection{Problem statement}\label{problem-statement}

For every fixed constant \(c>0\), is there a deterministic algorithm
that, given three explicitly represented \(n\times n\) matrices
\(A,B,C\), decides exactly whether

\[
AB=C
\]

in \(\widetilde O(n^2)\) bit operations in each of the following input
domains?

\begin{enumerate}
\def\labelenumi{\arabic{enumi}.}
\tightlist
\item
  Integer matrices with every input entry in \([-n^c,n^c]\).
\item
  Matrices over a finite field \(\mathbb F_q\) with \(q\le n^c\).
\end{enumerate}

Here \(\widetilde O(n^2)\) means \(O(n^2(\log n)^d)\) for a constant
\(d\) independent of the matrices and \(n\); constants may depend on
\(c\). Integers use binary encoding. For the finite-field case, use a
standard explicit field representation: \(q=p^r\), elements represented
by polynomials of degree less than \(r\) over \(\mathbb F_p\), and an
irreducible degree-\(r\) modulus supplied with the input. The field
representation is promised valid. Arithmetic and intermediate bit
lengths contribute to the running time.

The answer must be correct on every input. There is no sparsity, rank,
symmetry, or nonsingularity promise, and no auxiliary certificate from a
prover. This is one derandomization question across the two domains
explicitly identified in the source. The field-encoding convention makes
its bit-cost interpretation explicit.

Verification is a basic task for checking exact matrix computations.
Freivalds's randomized test already has near-quadratic bit complexity
for these domains. Counting operations on arbitrarily large integers at
unit cost would change the question.

\newpage
\subsection{References}\label{references}

Huck Bennett, Karthik Gajulapalli, Alexander Golovnev, and Evelyn
Warton, \href{https://arxiv.org/pdf/2309.16176v2}{\emph{Matrix
Multiplication Verification Using Coding Theory}}, RANDOM 2024,
\href{https://doi.org/10.4230/LIPIcs.APPROX/RANDOM.2024.42}{LIPIcs 317,
Article 42}; full version §1.3, p.~10, and Definition 2.6, p.~12.
Section 1.2 discusses why an existing quadratic arithmetic-operation
algorithm is not a quadratic bit algorithm.

The same authors,
\href{https://arxiv.org/pdf/2508.10250}{\emph{Output-Sparse Matrix
Multiplication Using Compressed Sensing}}, 2025 preprint, accepted at
RANDOM 2026; §1.1, Theorems 1--2, and §1.2 explain the later
sparse-product advances and their relationship to verification.

\subsection{Status check --- 2026-09-10}\label{status-check-2026-09-10}

Checked Bennett--Gajulapalli--Golovnev--Warton v2, §§1.1--1.3, the 2025
follow-up, and targeted deterministic-verification/2025--2026 searches.
The 2024 paper recounts a deterministic near-quadratic algorithm under
the promise that \(AB-C\) has at most \(n\) nonzero entries, a
substantive solved input subclass. Its newer coding-theory and
sparse-product results do not yield the required unrestricted
near-quadratic bit bound. No such general algorithm was located. The
paper's lower-bound barrier concerns a restricted class of
linear-algebraic algorithms; it is not an unconditional impossibility
theorem. Unbounded unit-cost integer arithmetic changes the target
model.
\end{document}
